\section{设计原则提炼}

从 Anthropic 的实践中，我们可以提炼出以下适用于所有长时间运行 Agent 系统的设计原则：

\subsection{原则一：状态外部化}

Agent 的工作状态不应仅存在于 context window 中，而应该被外部化到持久存储（文件系统、数据库等）中。这样即使 context window 被清空，新的 session 也能从外部状态中恢复。

\subsection{原则二：任务分解与增量推进}

复杂任务必须被分解为可以在单个 session 内完成的小任务。每个 session 的目标应该是"完成一个功能并留下干净的状态"，而非"尽可能多地完成任务"。

\subsection{原则三：不可变的评判标准}

任务完成的评判标准（如 feature list）应该在一开始就确定，并且在执行过程中不可被 Agent 修改。这防止了 Agent 通过降低标准来"完成"任务。

\subsection{原则四：先验证后推进}

每个 session 开始时应该先验证现有功能的正确性，再开始新功能的开发。这防止了在已有问题的基础上叠加更多问题。

\subsection{原则五：工具即能力边界}

Agent 的能力上限由其可用工具决定。要提升 Agent 的表现，首先应该考虑是否可以提供更好的工具，而不仅仅是优化 prompt。

\begin{knowledgebox}{与人类工程实践的对应}
这些原则与人类软件工程的最佳实践高度一致：状态外部化对应文档和 wiki，任务分解对应 sprint planning，不可变标准对应 acceptance criteria，先验证后推进对应 CI/CD，工具即能力对应开发环境。这并非巧合——Anthropic 明确表示他们的灵感来自"优秀的软件工程师每天都在做的事情"。
\end{knowledgebox}

\subsection{本章小结}

从 Anthropic 的实践中可以提炼出五条通用设计原则：状态外部化、任务分解与增量推进、不可变的评判标准、先验证后推进、工具即能力边界。这些原则不仅适用于编码 Agent，也可以推广到其他类型的长时间运行 Agent 系统。


